\section{Data}
\label{sec:data}

\subsection{Data sources}
\label{sec:ft-sources}

\begin{table}[h]
\centering
\small
\begin{tabular}{llll}
\hline
Source & Purpose & Verified labels & Used for \\
\hline
HierSVA~\cite{hiersva2026} & Hierarchical contracts & Yes & SFT/RL \\
OpenTitan~\cite{opentitan} & Assume-guarantee structure & Yes & SFT/RL \\
HWSec-UNC~\cite{ahmad2024hwsec} & Security properties & Yes & SFT/RL \\
riscv-formal~\cite{riscvformal} & Compositional verification & Yes & SFT/RL \\
VERT~\cite{menon2025vert} & SVA syntax breadth & Mixed & SFT \\
Verixa~\cite{verixa2025} & Verification context & Mixed & SFT \\
NVIDIA CVDP~\cite{nvidiacvdp2025} & Specification and RTL & Mixed & SFT \\
VerilogEval~\cite{liu2023verilogeval} & English and RTL & No & SFT \\
\hline
\end{tabular}
\caption{Fine-tuning and reinforcement-learning data sources.}
\label{tab:ft-sources}
\end{table}

\subsection{Data augmentation}
\label{sec:ft-augmentation}

We augment each verified pair of a specification and contract in two ways.
For identifier renaming, we replace signal, port, and module names with sampled alternatives that follow the design's naming conventions.
For syntactic paraphrasing, we apply rewrite rules to generate semantically equivalent SystemVerilog and SVA variants by reordering conjuncts or using De Morgan's laws and equivalent temporal operators.

\subsection{Designs and assertions for the environment}
\label{sec:rl-instances}

Each instance consists of a design and a target assertion. We include only
assertions known to hold, which lets us attribute failures to the agent.

We use design-level Training, Validation, and Testing partitions.
The partitions are the same for fine-tuning and RL.
Every assertion, generated variant, and proof trace derived from a design stays in that design's partition.
